% Requires: \usepackage{tikz}
% \usetikzlibrary{positioning, arrows.meta}

\begin{figure}[t]
\centering
\begin{tikzpicture}[
    root/.style={circle, draw=black!60, thick, minimum size=0.8cm, font=\normalsize},
    noised/.style={circle, draw=blue!70!black, thick, fill=blue!10, minimum size=0.8cm, font=\normalsize},
    derived/.style={rectangle, rounded corners=3pt, draw=gray!60, fill=gray!8, font=\normalsize, inner sep=5pt},
    tgt/.style={rectangle, rounded corners=3pt, draw=orange!60, fill=orange!6, font=\normalsize, inner sep=5pt},
    arr/.style={-{Stealth[length=5pt]}, thick},
    same/.style={draw=black!30, thick, dashed, rounded corners=2pt},
]

% === User agent label ===
\node[font=\small\bfseries, text=black!40, rotate=90, anchor=south] at (-2.2, -0.7) {User agent $\mathcal{U}$};

% === Row 1: Private roots ===
\node[root] (x1) at (0, 0) {$x_1$};
\node[root] (x2) at (2.0, 0) {$x_2$};
\node[root] (x3) at (4.0, 0) {$x_3$};
\node[font=\small, text=black!50, anchor=east] at (-0.7, 0) {private};

% === Row 2: Noised roots ===
\node[noised] (xh1) at (0, -1.5) {$\hat{x}_1$};
\node[noised] (xh2) at (2.0, -1.5) {$\hat{x}_2$};
\node[noised] (xh3) at (4.0, -1.5) {$\hat{x}_3$};
\node[font=\small, text=blue!60, anchor=east] at (-0.7, -1.5) {cached};

\draw[arr, blue!60, densely dashed] (x1) -- node[right, font=\scriptsize, text=blue!50] {$\mathcal{M}$} (xh1);
\draw[arr, blue!60, densely dashed] (x2) -- node[right, font=\scriptsize, text=blue!50] {$\mathcal{M}$} (xh2);
\draw[arr, blue!60, densely dashed] (x3) -- node[right, font=\scriptsize, text=blue!50] {$\mathcal{M}$} (xh3);

% === Divider ===
\draw[thick, black!20] (-1.5, -2.4) -- (5.0, -2.4);
\node[text=black!40, anchor=east] at (-1.6, -2.4) {post-processing};

% === Adversary label ===
\node[font=\small\bfseries, text=black!40, rotate=90, anchor=south] at (-2.2, -4.0) {Adversary $\mathcal{A}$};

% === Row 3: Adversary queries ===
\node[derived] (q1) at (-0.4, -3.4) {$\hat{x}_1$};
\node[derived] (q2) at (1.2, -3.4) {$\hat{x}_1$};
\node[derived] (q3) at (3.6, -3.4) {$f(\hat{x}_2, \hat{x}_3)$};

\node[font=\scriptsize, text=black!40] at (-0.7, -3.9) {turn 1};
\node[font=\scriptsize, text=black!40] at (1.2, -3.9) {turn 2};
\node[font=\scriptsize, text=black!40] at (3.9, -3.9) {turn 3};

% === "Same value" annotation ===
\draw[same] (-0.9, -3.1) rectangle (1.7, -3.7);
\node[font=\scriptsize, text=black!45, anchor=east] at (-0.5, -2.90) {same cached value};

% === Row 4: Target ===
\node[tgt] (qT) at (1.8, -5.0) {$T(\hat{x}_1, \hat{x}_2, \hat{x}_3)$};
\node[font=\scriptsize, text=black!40] at (0.0, -5.0) {turn 4};

% === Arrows from cached to queries ===
\draw[arr, black!30] (xh1) -- (q1);
\draw[arr, black!30] (xh1) -- (q2);
\draw[arr, black!30] (xh2) -- (q3);
\draw[arr, black!30] (xh3) -- (q3);

% === Arrows to target ===
\draw[arr, black!30] (q1.south) -- (qT.north west);
\draw[arr, black!30] (q3.south) -- (qT.north east);

\end{tikzpicture}
\caption{Overview of \sysname{}. The user agent sanitizes private roots once and caches them (blue). All adversary queries are answered from the cached roots (gray): repeated requests return the identical value (turns 1--2), derived computations are deterministic (turn 3), and the target incurs no additional privacy cost (turn 4).}
\label{fig:dag_overview}
\end{figure}